\documentclass[11pt]{article}

% Packages
\usepackage{kotex}
\usepackage{graphicx}
\usepackage{indentfirst}
\usepackage{setspace}
\usepackage{amssymb}
\usepackage{pifont}
\usepackage[linguistics]{forest}       % Syntax Tree 패키지 
\usepackage[a4paper, 
            top = 2.5cm, 
            bottom = 2.5cm, 
            left = 2cm, 
            right = 2cm]{geometry}     % 용지 여백 조절 패키지 
\usepackage{fancyhdr}                  % 머리글, 바닥글 조정 패키지
\usepackage{enumitem}
\pagestyle{fancy}
\fancyhf{} % 기존 머리글과 바닥글 초기화
\fancyhead[R]{언어습득: quiz \#1}   % 오른쪽: 문서 제목
\fancyfoot[L]{\nouppercase{\leftmark}} % 왼쪽: 현재 section 제목
\fancyfoot[R]{\thepage}                % 오른쪽: 페이지 번호
\renewcommand{\sectionmark}[1]{%
  \markboth{#1}{}%
}
\renewcommand{\headrulewidth}{0.4pt}   % 머리글 아래쪽 선
\renewcommand{\footrulewidth}{0.4pt}   % 바닥글 위쪽 선
\setlength{\headheight}{14pt}          % 머리글 높이 확보

% 제목 페이지용 바닥글
\fancypagestyle{firstpage}{%
  \fancyhf{}%
  \fancyfoot[L]{1. Tree}%
  \fancyfoot[R]{\thepage}%
  \renewcommand{\headrulewidth}{0pt}%
  \renewcommand{\footrulewidth}{0.4pt}%
}

\setcounter{secnumdepth}{0}            % section 넘버링 제거

% 동그라미 숫자의 오른쪽 위에 체크 표시를 겹쳐 출력
\newcommand{\checkedcircled}[1]{%
  \tikz[baseline=(number.base)]{%
    \node[inner sep=0pt] (number) {\textcircled{#1}};
    \draw[overlay, line width=1.1pt, line cap=round, line join=round]
      ([xshift=-0.45em,yshift=0.39em]number.center) --
      ([xshift=0em,yshift=-0.08em]number.center) --
      ([xshift=0.47em,yshift=0.38em]number.center);
  }%
}

\forestset{
  terminal/.style={
    no edge,
    l=3ex,
    for parent={l sep=0pt}
  }
}

% Short form for a lexical category and its terminal.
\forestset{
  lex/.style={append={[#1, terminal]}},
  simple terminals/.style={
    delay={for tree={split option={content}{ }{content,lex}}}
  }
}

% Reusable GPS tree block; default width is 0.48\textwidth.
\NewDocumentEnvironment{gpstree}{O{0.48\textwidth} m +b}
{%
  \begin{minipage}[t]{#1}
  \vspace{0pt}%
  \textbf{#2}\par
  \medskip
  \centering
  \makebox[\linewidth][c]{%
    \begin{forest} simple terminals
      #3
    \end{forest}%
  }
  \end{minipage}%
}
{}

% Metadata
\title{2026 Fall 언어습득: \\ \medskip \large{quiz \#1}}
\author{황효성 / 83-604}
\date{October 12, 2026}

\begin{document}
\maketitle

\subsection{1. Skinner의 설명에 따르면, 누군가 그림을 보고 "Dutch"라고 말하는 것은 그림의 \\ '네덜란드풍'이라는 속성이 그 반응을 통제했기 때문이다. Chomsky(1959)가 이런 설명을 비판한 핵심은 무엇인가?}
\medskip
\noindent
\begin{enumerate}[label=\textcircled{\arabic*}]
    \item 외부에서 관찰할 수 있는 자극을 언어 설명에 사용했다는 점이다. 
    \item 발화가 나타나기 전에 아동의 모방 능력을 검사하지 않았다는 점이다. 
    \item[\checkedcircled{3}]\setcounter{enumi}{3}%
    발화 결과에 맞추어 자극을 정해, 그 자극으로 발화를 예측하기 어렵다는 점이다. 
    \item 한 그림에서 서로 다른 속성에 주의를 기울일 수 있다고 본다는 점이다. 
\end{enumerate}
\bigskip

\subsection{2. Mayberry et al.(2002)에서 영어 학습 시작 연령이 비슷한 두 선천성 놓인 집단 중 영아기에 수어(ASL)를 경험한 집단이 초기 언어 경험이 거의 없는 집단보다 영어 문법성 판단 수행이 높았다. 이 비교에 대한 해석으로 가장 적절한 것은?}
\medskip
\noindent
\begin{enumerate}[label=\textcircled{\arabic*}]
    \item 농인의 낮은 영어 수행은 선천성 청각 손실만으로 설명할 수 있다. 
    \item[\checkedcircled{2}] 
    \setcounter{enumi}{2}%
    선청성 농인 간 수행 차이를 설명할 때 초기 언어 경험을 고려해야 한다. 
    \item 두 집단의 수행 차이는 초기 언어 경험 하나만의 인과적 효과로 확정된다. 
    \item 초기 ASL 경험은 이후 수어 학습을 돕지만 영어 학습과는 관련이 없다. 
\end{enumerate}
\bigskip

\subsection{3. Tomasello \& Farrar(1986)의 새 단어 학습 실험에서, 아동이 이미 몰두한 사물의 이름을 알려준 따라가기 조건은 아무 사물에도 몰두하지 않을 때 사물을 들어 보이며 알려준 지시 조건과 어떻게 달랐는가?}
\medskip
\noindent
\begin{enumerate}[label=\textcircled{\arabic*}]
    \item[\checkedcircled{1}] 
    \setcounter{enumi}{1}% 
    단어 이해가 더 높았지만, 산출에서는 뚜렷한 조건 차이를 확인하지 못했다. 
    \item 단어 산출이 더 높았지만, 이해에서는 뚜렷한 조건 차이를 확인하지 못했다. 
    \item 단어 이해와 산출 모두에서 두 조건의 결과가 같았다. 
    \item 단어 이해와 산출 모두에서 지시 조건의 결과가 더 높았다. 
\end{enumerate}
\bigskip

\subsection{4. Fernald et al.(2008)의 LWL 절차에서, 한 아동이 어떤 시행에서 표적어가 시작될 때 이미 표적 그림을 보고 있었고 분석 구간 내내 그 그림을 보았다. 이 시행은 정확도 분석과 반응 시간 분석에서 각각 어떻게 처리하는가?}
\medskip
\noindent
\begin{enumerate}[label=\textcircled{\arabic*}]
    \item 정확도에 포함하고, 반응 시간은 0ms로 계산한다. 
    \item 정확도에서 제외하고, 반응 시간은 0ms로 계산한다. 
    \item 정확도에서 제외하고, 반응 시간 분석에서도 제외한다. 
    \item[\checkedcircled{4}] 
    \setcounter{enumi}{4}%
    정확도에 포함하고, 반응 시간 분석에서는 제외한다. 
\end{enumerate}
\bigskip

\subsection{5. Yang(2004)의 'UG가 인도하는 확률학습'에서 보편문법(UG)과 통계학습의 역할을 가장 적절하게 연결한 것은? }
\medskip
\noindent
\begin{enumerate}[label=\textcircled{\arabic*}]
    \item UG가 입력의 빈도를 계산하고, 통계학습이 가능한 문법의 범위를 정한다. 
    \item[\checkedcircled{2}] 
    \setcounter{enumi}{2}% 
    UG가 가능한 문법을 제한하고, 통계학습이 입력에 따라 후보의 가중치를 조정한다. 
    \item UG가 개별 언어의 문법을 완성하고, 통계학습이 발음의 개인차를 조정한다. 
    \item UG가 학습 결과를 일반화하고, 통계학습이 언어에 고유한 선천적 제약을 만든다.  
\end{enumerate}
\bigskip

\subsection{6. 님 침스키 프로젝트에서 신호를 여러 개 조합했다는 사실만으로 인간 아동과 같은 통사 발달을 인정하기 어려웠던 근거는?}
\medskip
\noindent
\begin{enumerate}[label=\textcircled{\arabic*}]
    \item 학습한 신호 수가 늘었지만, 의사소통의 수단으로 음성 대신 수어를 사용했다. 
    \item 학습한 신호가 100개에 미치지 못해 조합에 사용할 신호가 부족했다. 
    \item[\checkedcircled{3}] 
    \setcounter{enumi}{3}% 
    조합에 반복이 많고 평균 발화 길이가 늘지 않았으며, 교사가 직전에 한 신호에 의존했다. 
    \item 배운 신호들을 새로운 상황에 조합했고, 교사의 신호 없이 자발적으로 사용했다. 
\end{enumerate}
\bigskip

\subsection{7. Hartshorne et al.(2018)은 문법 학습 능력이 약 17세까지 비교적 유지되지만, 원어민 수준에 도달하려면 약 10--12세까지는 학습을 시작해야 한다고 추정했다. 두 시점이 다른 이유로 가장 적절한 것은?}
\medskip
\noindent
\begin{enumerate}[label=\textcircled{\arabic*}]
    \item[\checkedcircled{1}] 
    \setcounter{enumi}{1}% 
    최종 수준은 학습 능력의 변화와 숙달에 필요한 시간에 함께 좌우되기 때문이다. 
    \item 시작 연령이 달라도 같은 기간 동안 배우면 최종 수준은 같아지기 때문이다. 
    \item 학습 능력이 감소하는 시점이 곧 원어민 수준에 도달할 수 있는 마지막 학습 시작 시점이기 때문이다. 
    \item 문법 학습 능력이 감소하기 시작하면 이후의 언어 경험은 학습에 기여하지 않기 때문이다.  
\end{enumerate}
\bigskip

\subsection{8. O'Doherty et al.(2011)에서 약 2세 아동은 녹화 영상 속 어른이 자신에게 새 사물 이름을 말하는 조건에서는 학습하지 못했지만, 영상 속 두 어른이 그 사물을 주고받으며 이름을 말하는 장면을 본 조건에서는 학습했다. 이 결과가 지지하는 해석은?}
\medskip
\noindent
\begin{enumerate}[label=\textcircled{\arabic*}]
    \item 영상 속 화자의 직접 호명은 실제 대화에서 반응을 주고받는 경험을 대신한다. 
    \item 직접 참여한 대화에서 먼저 배운 단어가 다른 사람의 대화를 이해하는 데 전이되었다. 
    \item 화면의 제약을 보완하려면 녹화 영상을 실시간 영상 통화로 바꾸어야 한다. 
    \item[\checkedcircled{4}] 
    \setcounter{enumi}{4}% 
    타인들이 상호작용하는 장면을 관찰하는 것도 단어 학습에 도움이 될 수 있다.   
\end{enumerate}
\bigskip

\subsection{9. 한 아동의 자연발화를 한 달 간격으로 30분씩 녹음했다. 특정 문법 형태가 첫 녹음에서는 나오지 않았고 두 번째 녹음에서 처음 관찰되었다. 이 자료로 내릴 수 있는 결론은?}
\medskip
\noindent
\begin{enumerate}[label=\textcircled{\arabic*}]
    \item 두 녹음 사이에 그 형태를 습득했으며, 첫 녹음은 미습득 상태를 보여 준다. 
    \item 두 번째 녹음일이 그 형태의 습득일이며, 그 전에 사용했을 가능성은 낮다. 
    \item[\checkedcircled{3}] 
    \setcounter{enumi}{3}% 
    표본에서 처음 관찰된 시점은 알지만, 실제 습득 시점은 확정하기 어렵다. 
    \item 녹음 간격이 일정하므로, 해당 형태의 발달이 불연속적이라고 판단할 수 있다.   
\end{enumerate}
\bigskip

\subsection{10. 고진폭 빨기(HAS) 실험에서 [ba]가 반복되자 영아의 빨기 반응이 감소했다. [pa]로 바꾼 집단에서는 반응이 다시 증가했지만, [ba]를 계속 들려준 집단에서는 낮게 유지되었다. 이 결과가 직접 뒷받침하는 것은?}
\medskip
\noindent
\begin{enumerate}[label=\textcircled{\arabic*}]
    \item 두 음절 중 [pa]가 [ba]보다 영아에게 더 선호된다. 
    \item[\checkedcircled{2}] 
    \setcounter{enumi}{2}% 
    영아는 [ba]와 [pa]의 차이를 지각할 수 있다. 
    \item 영아는 두 음절로 구별되는 단어의 뜻을 이해한다. 
    \item 영아는 들은 [pa]를 자신의 발음으로 산출할 수 있다. 
\end{enumerate}
\bigskip

\end{document}
